% Generated by scripts/assemble.py; corrections belong in source/corrections.json.
\examhead{125}{中国海洋大学 2025 年数学分析试卷(回忆版)}{119}
% p120-b09; source PDF page 120
\par\Needspace{5\baselineskip}
\blockanchor{p120-b09}
\majorquestion{一}计算题.\par
% p120-b10; source PDF page 120
\blockanchor{p120-b10}
\question{1}计算 $\lim_{n\to \infty}n^2 (\sqrt[n]{2} -\sqrt[n + 1]{2}).$\par
% p120-b11; source PDF page 120
\blockanchor{p120-b11}
\question{2}忘了.\par
% p120-b12; source PDF page 120
\blockanchor{p120-b12}
\question{3}计算 $\int_0^{+\infty} \frac{e^{2x} - e^x}{x} \, \mathrm{d}x$. (需判断可计算性)\par
% p120-b13; source PDF page 120
\par\Needspace{5\baselineskip}
\blockanchor{p120-b13}
\majorquestion{二}判断题.\par
% p120-b14; source PDF page 120
\blockanchor{p120-b14}
\question{1}设函数列 $\{f_n(x)\}, \{g_n(x)\}$ 分别一致收敛于 $f(x), g(x)$, 则 $\{f_n(x)g_n(x)\}$ 是否一致收敛于 $f(x)g(x)$?
\question{2}$a_n > 0, \{a_n\}$ 单调递减, $\sum_{n=0}^{\infty} a_n$ 收敛, 则 $\lim_{n \to \infty} na_n = 0$.\par
% p120-b15; source PDF page 120
\blockanchor{p120-b15}
\question{3}判断积分 $\int_{2}^{+\infty}\frac{\sin^2x}{x^\lambda(x^\lambda + \sin x)}\mathrm{d}x$ 的敛散性.\par
% p120-b16; source PDF page 120
\par\Needspace{5\baselineskip}
\blockanchor{p120-b16}
\majorquestion{三}证明题.\par
% p120-b17; source PDF page 120
\blockanchor{p120-b17}
\question{1}设 $f(x, y)$ 在 $(x_0, y_0)$ 连续, $f_x, f_y$ 可微, 证明: $f_{xy}(x_0, y_0) = f_{yx}(x_0, y_0)$.\par
% p120-b18; source PDF page 120
\par\Needspace{6\baselineskip}
\blockanchor{p120-b18}
\question{2}设正项级数 $\sum_{n=1}^{\infty} u_n$ 收敛, 则对其进行适当重新排列组合得到的新数列\par
\nopagebreak[4]
% p120-b19; source PDF page 120
\blockanchor{p120-b19}
\displayformula{(u _ {k 1}, u _ {k 2}, \dots , u _ {k m}), \dots , (u _ {k ^ {\prime} 1}, u _ {k ^ {\prime} 2}, \dots , u _ {k ^ {\prime} l}), \dots}
% p120-b20; source PDF page 120
\blockanchor{p120-b20}
求和也收敛.\par
% p120-b21; source PDF page 120
\par\Needspace{5\baselineskip}
\blockanchor{p120-b21}
\majorquestion{四}隐函数 $F(x, y) = x^3 + y^3 - 3xy$.\par
% p120-b22; source PDF page 120
\blockanchor{p120-b22}
\question{1}求隐函数的渐进线.\par
% p120-b23; source PDF page 120
\blockanchor{p120-b23}
\question{2}求  $y = f(x)$  的二阶导数和极值.\par
% p120-b24; source PDF page 120
\par\Needspace{5\baselineskip}
\blockanchor{p120-b24}
\majorquestion{五}用有限覆盖定理证明闭区间上的零点存在定理.\par
% p120-b25; source PDF page 120
\par\Needspace{5\baselineskip}
\par\Needspace{6\baselineskip}
\blockanchor{p120-b25}
\majorquestion{六}$D = \{(x,y) \mid x^2 + y^2 \leq 1\}$, 设 $H(x,y)$ 是 $D$ 上的调和函数, 即 $H_{xx} + H_{yy} = 0$. 证明: 若\par
\nopagebreak[4]
% p120-b26; source PDF page 120
\blockanchor{p120-b26}
\displayformula{H (x, y) \equiv 1, (x, y) \in \partial D.}
% p120-b27; source PDF page 120
\blockanchor{p120-b27}
则 $H(x,y)\equiv 1,(x,y)\in \overline{D}.$\par
% p121-b00; source PDF page 121
\par\Needspace{5\baselineskip}
\blockanchor{p121-b00}
\majorquestion{七}计算积分 $I = \iint_{S} \frac{x \, \mathrm{d}y \, \mathrm{d}z + y \, \mathrm{d}z \, \mathrm{d}x + z \, \mathrm{d}x \, \mathrm{d}y}{(x^2 + y^2 + z^2)^{\frac{3}{2}}}$, 其中 $S$ 是上半球面 $x^2 + y^2 + z^2 = 1$, 取上侧.
\majorquestion{八}求 $x^{2} + 2y^{2} + 3z^{2} = 1$ 与坐标平面 $xoy$ 围成区域的体积.\par
